\section{Introduction}
  \label{sec:introduction}

  The experimental violation of Bell inequalities stands as one of the most robust
  and conceptually challenging features of quantum theory.
  Repeated tests have confirmed that correlations observed between spacelike separated
  measurement outcomes cannot be reproduced by any local hidden-variable model
  satisfying statistical factorization.
  This empirical fact is commonly interpreted as requiring either nonlocal dynamics,
  the abandonment of realism, or more radical departures from classical causal intuition.

  Despite these conclusions, it is important to recall that Bell inequalities do not
  constrain physical dynamics directly.
  They constrain the logical and probabilistic structure of effective descriptions,
  specifically the assumption that joint outcome probabilities factorize into independent
  local contributions conditioned on a complete specification of relevant variables.
  The failure of Bell inequalities therefore indicates a breakdown of effective
  factorization, but does not, by itself, mandate superluminal influences or causal
  nonlocality.

  In this work, we explore a general and minimal mechanism by which Bell-type
  correlations may arise without invoking nonlocal dynamics, retrocausality, or hidden
  variables.
  The central observation is that effective physical descriptions typically arise through
  many-to-one mappings from underlying configurations to observable outcomes.
  When this mapping is non-injective, distinct underlying configurations correspond
  to the same observable description, leading generically to a loss of factorization at
  the level of observable probabilities.

  We further argue that such non-injectivity is not an arbitrary feature, but the natural
  consequence of a finite descriptive capacity of effective observables.
  Observable states correspond not to individual microscopic configurations, but to
  equivalence classes whose multiplicity reflects a bounded resolution of the effective
  description.
  In this sense, non-injectivity can be understood as the manifestation of a finite
  projective capacity, which limits the amount of distinguishable information that can
  be encoded in observable variables.

  Within this framework, joint probabilities arise from averaging over entire classes of
  underlying configurations.
  Even if the underlying dynamics is local and deterministic, the induced probability
  structure at the observable level does not admit a factorized representation.
  The violation of Bell inequalities is thus traced to a structural property of the
  descriptive mapping itself, rather than to any dynamical nonlocality.

  The effective loss of information associated with this projection is not arbitrary.
  In physical realizations, it is controlled by a finite resolution scale, set in quantum
  theory by Planck's constant \(h\), which determines the minimal distinguishability of
  underlying configurations.
  Quantum states therefore correspond to finite-capacity equivalence classes rather than
  sharply defined microscopic states.
  The breakdown of factorization reflects this irreducible coarse-graining, not an
  epistemic lack of knowledge.

  This perspective also provides a natural characterization of the quantum-to-classical
  transition.
  In regimes where the effective descriptive capacity is large compared to the occupied
  state space, the mapping becomes approximately injective.
  In this low-occupancy regime, observable probabilities recover an approximately
  factorized form and Bell violations are suppressed.
  Classical statistical independence thus appears as a limiting case of effective
  injectivity, rather than as a fundamental principle.

  The scope of the present work is deliberately restricted.
  We do not propose a modification of quantum dynamics, introduce additional physical
  degrees of freedom, or advance a specific underlying dynamical model.
  Our analysis is structural in nature and concerns the logical conditions under which
  Bell-type inequalities apply to effective descriptions.
  By isolating non-injective projection, understood as a finite-capacity effect, as a
  sufficient mechanism for the failure of factorization, we aim to clarify the conceptual
  origin of quantum correlations while preserving both operational locality and the
  empirical content of quantum theory.
  The present work does not introduce new physical predictions.
  Its aim is to identify structural constraints on the class of effective
  descriptions to which Bell-type inequalities apply.

  The non-injective projection considered here is not introduced ad hoc, but arises
  naturally in relational frameworks in which spacetime and observables emerge as
  effective descriptions from an underlying configuration space, as discussed
  in~\cite{Beau2026a}.
